\begin{textsample}{NEG Dim 1 – sp – Score 5.00 – sp/sp\_em\_7.txt}  \label{ex:f1_neg_125}
O compromisso com o desenvolvimento das competências O Currículo Paulista sinaliza a necessidade de que as decisões pedagógicas promovam o desenvolvimento de competências necessárias ao pleno desenvolvimento do estudante.
Reiterando os termos da BNCC, o Currículo Paulista define competência como “ a mobilização de conhecimentos ( conceitos e procedimentos ), habilidades ( práticas, cognitivas e socioemocionais ), atitudes e valores para resolver demandas complexas da vida cotidiana, do pleno exercício da cidadania e do mundo do trabalho ” ( BRASIL, 2018, p.8 ).
Assim, o Currículo indica claramente o que o estudante deve “ saber ” ( em termos de conhecimentos, habilidades, atitudes e valores ) e, sobretudo, o que deve “ saber fazer ”, considerando a mobilização desses conhecimentos, habilidades, atitudes e valores para resolver demandas complexas da vida cotidiana, do pleno exercício da cidadania e do mundo do trabalho.
Espera -se que essas indicações possam orientar as escolas para o fortalecimento de ações que assegurem ao estudante a transposição de conhecimentos, habilidades, atitudes e valores em intervenções concretas e solidárias ( aprender a fazer e a conviver ) no processo da construção de sua identidade, aprimorando as capacidades de situar -se e perceber -se na diversidade, de pensar e agir no mundo de modo empático, respeitoso à diversidade, criativo e crítico ( aprender a ser ), bem como de desenvolver sua autonomia para gerenciar a própria aprendizagem e continuar aprendendo ( aprender a aprender ).
É necessário garantir que, ao final da Educação Básica, o estudante paulista se constitua como cidadão autônomo, capaz de interagir de maneira crítica e solidária, de atuar de maneira consciente e eficaz nas ações que demandam análise criteriosa e na tomada de decisões que impactam o bem comum, de buscar e analisar criticamente diferentes informações e ter plena consciência de que a aprendizagem é demanda para a vida toda.

% matched lemmas: 
\end{textsample}
